
\documentclass{article}
\usepackage{colortbl}
\usepackage{makecell}
\usepackage{multirow}
\usepackage{supertabular}

\begin{document}

\newcounter{utterance}

\centering \large Interaction Transcript for game `eqbench', experiment `eq\_bench', episode 6 with reflexion{-}llama{-}3.1{-}8b{-}instruct.
\vspace{24pt}

{ \footnotesize  \setcounter{utterance}{1}
\setlength{\tabcolsep}{0pt}
\begin{supertabular}{c@{$\;$}|p{.15\linewidth}@{}p{.15\linewidth}p{.15\linewidth}p{.15\linewidth}p{.15\linewidth}p{.15\linewidth}}
    \# & \multicolumn{2}{c}{Player} && \multicolumn{2}{c}{Game Master} \\
    \hline

    \theutterance \stepcounter{utterance}  
    & & & \multicolumn{4}{p{0.6\linewidth}}{
        \cellcolor[rgb]{0.9,0.9,0.9}{
            \makecell[{{p{\linewidth}}}]{
                \texttt{\tiny{[P1$\langle$GM]}}
                \texttt{You have attempted this same task before and did not succeed. These are the plans you wrote after those attempts.} \\
\texttt{Trial \#1: New plan: } \\
\\ 
\texttt{Upon re{-}reading the dialogue, I realized that I focused too much on the surface{-}level conversation and didn't delve deeper into the emotional dynamics between Ellie and Gabe. I should have paid more attention to the underlying power struggle and the emotional cues that Gabe and Ellie exchange throughout the conversation.} \\
\\ 
\texttt{To improve my analysis, I will focus on identifying the emotional triggers and responses that Ellie and Gabe exhibit, and use those to inform my assessment of their emotional states. Specifically, I will look for:} \\
\\ 
\texttt{1. Verbal cues: Pay attention to the words and phrases that Ellie and Gabe use to express their emotions, such as "I don't lose, I learn" (Gabe's confidence) or "I have a masterpiece to create" (Ellie's competitiveness).} \\
\texttt{2. Nonverbal cues: Consider the tone, body language, and facial expressions that Ellie and Gabe display, such as Gabe's calm demeanor or Ellie's dismissive tone.} \\
\texttt{3. Emotional escalation: Identify the moments when the conversation escalates or de{-}escalates, and use those to gauge the emotional intensity of the characters.} \\
\\ 
\texttt{By incorporating these strategies, I will be able to provide a more nuanced and accurate assessment of Gabe's emotional state at the end of the dialogue.} \\
\\ 
\\ 
\\ 
\texttt{Your task is to predict the likely emotional responses of a character in this dialogue:} \\
\\ 
\texttt{Ellie: So, Gabe, are you ready to lose?} \\
\texttt{Gabe: I don't lose, Ellie. I learn.} \\
\texttt{Ellie: Ah, the eternal student. Must be nice living in a world where failure doesn't exist.} \\
\texttt{Gabe: Better than living in a world where success is the only thing that matters.} \\
\texttt{Ellie: And what's wrong with wanting to succeed? } \\
\texttt{Gabe: Nothing. But when it becomes an obsession, it blinds you.} \\
\texttt{Ellie: And who says I'm obsessed?} \\
\texttt{Gabe: You don't have to say it, Ellie. It's in your eyes, your work, your every breath.} \\
\texttt{Ellie: And what about you, Gabe? What's in your eyes? Fear? Doubt? Or just a whole lot of nothing?} \\
\texttt{Gabe: I guess you'll never know.} \\
\texttt{Ellie: I guess I won't. And maybe I don't want to.} \\
\texttt{Gabe: Maybe you don't. But maybe you do. Maybe you're just too scared to admit it.} \\
\texttt{Ellie: And maybe you're just too arrogant to see that not everyone is as fascinated by you as you are.} \\
\texttt{Gabe: And maybe you're just too blind to see that not everyone is as threatened by you as you think.} \\
\texttt{Ellie: You know what, Gabe? I don't have time for this. I have a masterpiece to create.} \\
\texttt{Gabe: And I have a student to teach. Good luck, Ellie. You're going to need it.} \\
\texttt{Ellie: I don't need luck, Gabe. I have talent.} \\
\texttt{{[}End dialogue{]}} \\
\\ 
\texttt{At the end of this dialogue, Gabe would feel...} \\
\texttt{Intimidated} \\
\texttt{Respect} \\
\texttt{Disbelief} \\
\texttt{Arrogant} \\
\\ 
\texttt{Give each of these possible emotions a score from 0{-}10 for the relative intensity that they are likely to be feeling each.} \\
\\ 
\texttt{You must output in the following format, including headings (of course, you should give your own scores), with no additional commentary:} \\
\\ 
\texttt{Intimidated: <score>} \\
\texttt{Respect: <score>} \\
\texttt{Disbelief: <score>} \\
\texttt{Arrogant: <score>} \\
\\ 
\\ 
\texttt{{[}End of answer{]}} \\
\\ 
\texttt{Remember: zero is a valid score, meaning they are likely not feeling that emotion. You must score at least one emotion > 0.} \\
\\ 
\texttt{Your answer:} \\
            }
        }
    }
     \\ \\

    \theutterance \stepcounter{utterance}  
    & \multicolumn{4}{p{0.6\linewidth}}{
        \cellcolor[rgb]{0.561,0.737,0.561}{
            \makecell[{{p{\linewidth}}}]{
                \texttt{\tiny{[P1$\rangle$GM]}}
                \texttt{Intimidated: 8} \\
\texttt{Respect: 0} \\
\texttt{Disbelief: 0} \\
\texttt{Arrogant: 6} \\
            }
        }
    }
    & & \\ \\

    \theutterance \stepcounter{utterance}  
    & & & \multicolumn{2}{p{0.3\linewidth}}{
        \cellcolor[rgb]{0.9,0.9,0.9}{
            \makecell[{{p{\linewidth}}}]{
                \texttt{\tiny{[GM$|$GM]}}
                \texttt{Intimidated: 8} \\
\texttt{Respect: 0} \\
\texttt{Disbelief: 0} \\
\texttt{Arrogant: 6} \\
            }
        }
    }
    & & \\ \\

    \theutterance \stepcounter{utterance}  
    & & & \multicolumn{2}{p{0.3\linewidth}}{
        \cellcolor[rgb]{0.9,0.9,0.9}{
            \makecell[{{p{\linewidth}}}]{
                \texttt{\tiny{[GM$|$GM]}}
                \texttt{Intimidated: 0} \\
\texttt{Respect: 0} \\
\texttt{Disbelief: 3} \\
\texttt{Arrogant: 4} \\
            }
        }
    }
    & & \\ \\

    \theutterance \stepcounter{utterance}  
    & & & \multicolumn{2}{p{0.3\linewidth}}{
        \cellcolor[rgb]{0.9,0.9,0.9}{
            \makecell[{{p{\linewidth}}}]{
                \texttt{\tiny{[GM$|$GM]}}
                \texttt{game\_result = LOSE} \\
            }
        }
    }
    & & \\ \\

    \theutterance \stepcounter{utterance}  
    & & \multicolumn{4}{p{0.6\linewidth}}{
        \cellcolor[rgb]{0.95,0.95,0.95}{
            \makecell[{{p{\linewidth}}}]{
                \texttt{\tiny{[P1$\leftrightarrow$P]}}
                \texttt{Post{-}episode Reflexion (not sent to GM):} \\
\texttt{New plan: } \\
\\ 
\texttt{Upon re{-}reading the dialogue, I realized that I focused too much on the surface{-}level conversation and didn't delve deeper into the emotional dynamics between Ellie and Gabe. I should have paid more attention to the underlying power struggle and the emotional cues that Gabe and Ellie exchange throughout the conversation.} \\
\\ 
\texttt{To improve my analysis, I will focus on identifying the emotional triggers and responses that Ellie and Gabe exhibit, and use those to inform my assessment of their emotional states. Specifically, I will look for:} \\
\\ 
\texttt{1. Verbal cues: Pay attention to the words and phrases that Ellie and Gabe use to express their emotions, such as "I don't lose, I learn" (Gabe's confidence) or "I have a masterpiece to create" (Ellie's competitiveness).} \\
\texttt{2. Nonverbal cues: Consider the tone, body language, and facial expressions that Ellie and Gabe display, such as Gabe's calm demeanor or Ellie's dismissive tone.} \\
\texttt{3. Emotional escalation: Identify the moments when the conversation escalates or de{-}escalates, and use those to gauge the emotional intensity of the characters.} \\
\\ 
\texttt{To specifically address the emotions I scored in my previous attempt, I will:} \\
\\ 
\texttt{* Increase the score for "Intimidated" to 9, as Gabe's calm and confident demeanor may be perceived as intimidating by Ellie, especially when he says "I don't lose, I learn" and "I have a student to teach."} \\
\texttt{* Decrease the score for "Respect" to 2, as the conversation between Ellie and Gabe is more confrontational than respectful.} \\
\texttt{* Decrease the score for "Disbelief" to 1, as both characters seem to be genuinely expressing their emotions and opinions.} \\
\texttt{* Increase the score for "Arrogant" to 8, as Ellie's dismissive tone and competitive language ("I have a masterpiece to create") suggest a sense of arrogance and superiority.} \\
            }
        }
    }
    & \\ \\

\end{supertabular}
}

\end{document}
